% ============================================================
% Section 54: 一维双原子链的热容——M≫m极限
% ============================================================
\newpage
\section{固体理论：一维双原子链的热容——$M\gg m$极限}
\label{sec:1d-diatomic-heat-capacity}

\begin{modelbox}[题目：一维复式格子的晶格热容量]
考虑由两种不同原子（质量分别为 $M$ 和 $m$）所构成的一维复式格子，相邻同种原子间距离为 $2a$，相互作用恢复力常数为 $\beta$。设 $M\gg m$，求晶格热容量，并讨论高低温极限情况。
\end{modelbox}

\begin{tipbox}[考点快速分析]
\begin{itemize}
    \item \textbf{一维双原子链色散}：两支格波——声学支（同相振动）和光学支（反相振动）
    \item \textbf{$M\gg m$ 近似}：光学支频率几乎与 $q$ 无关（爱因斯坦模型）；声学支等效为轻原子静止、重原子振动的单原子链
    \item \textbf{一维模式密度}：$D(\omega)\propto 1/\sqrt{\omega_0^2-\omega^2}$（范霍夫奇点）
    \item \textbf{高温极限}：杜隆-珀蒂定律，每原子 $k_B$，共 $2N$ 个原子 $\Rightarrow C_V\to2Nk_B$
    \item \textbf{低温极限}：光学支指数冻结，声学支主导，$C_V\propto T$（一维特征，区别于三维德拜 $T^3$ 律）
\end{itemize}
\end{tipbox}

\begin{derivationbox}[标准解题过程]
\textbf{Step 0: 色散关系的近似处理}

一维双原子链的色散关系由运动方程导出。设试探解为
\begin{equation}
x_{2n+1}=Ae^{i[q(2n+1)a-\omega t]},\qquad x_{2n+2}=Be^{i[q(2n+2)a-\omega t]},
\end{equation}
代入运动方程后得到 $A,B$ 的齐次方程组，由系数行列式为零给出色散关系：
\begin{equation}
mM\omega^4-2\beta(m+M)\omega^2+4\beta^2\sin^2(qa)=0.
\end{equation}

解得两支格波：
\begin{equation}
\omega_{\pm}^2=\frac{\beta}{mM}\Bigl[(m+M)\pm\sqrt{(m+M)^2-4mM\sin^2(qa)}\Bigr].
\end{equation}

\textbf{光学支}（$\omega_+$）：在 $M\gg m$ 时，$m/M\ll1$，忽略 $q$ 依赖的高阶小量：
\begin{equation}
\omega_+^2\approx\frac{2\beta}{m}\Bigl(1+\frac{m}{M}\Bigr)\approx\frac{2\beta}{m}\equiv\omega_0^2.
\end{equation}
光学支近似为与 $q$ 无关的常数，可用\textbf{爱因斯坦模型}处理。

\textbf{声学支}（$\omega_-$）：在 $M\gg m$ 和长波极限下：
\begin{equation}
\omega_-^2\approx\frac{2\beta}{M}\sin^2(qa)\equiv\omega_0'^2\sin^2(qa),
\qquad \omega_0'=\sqrt{\frac{2\beta}{M}}.
\end{equation}
这正是轻原子 $m$ 静止、重原子 $M$ 构成的有效单原子链的色散关系。

\textbf{Step 1: 光学支对热容量的贡献（爱因斯坦模型）}

光学支共 $N$ 个模式（$N$ 为原胞数），每个模式频率均为 $\omega_0$。总振动能：
\begin{equation}
E_+=\sum_{\text{光学支}}\frac{\hbar\omega_0}{e^{\hbar\omega_0/k_BT}-1}
=\frac{N\hbar\omega_0}{e^{\hbar\omega_0/k_BT}-1}.
\end{equation}

热容量：
\begin{equation}
C_V^+=\frac{\partial E_+}{\partial T}
=Nk_B\Bigl(\frac{\hbar\omega_0}{k_BT}\Bigr)^2\frac{e^{\hbar\omega_0/k_BT}}{\bigl(e^{\hbar\omega_0/k_BT}-1\bigr)^2}
=Nk_B\Bigl(\frac{\Theta_E}{T}\Bigr)^2\frac{e^{\Theta_E/T}}{\bigl(e^{\Theta_E/T}-1\bigr)^2},
\end{equation}
其中爱因斯坦温度 $\Theta_E=\hbar\omega_0/k_B$。

\textbf{Step 2: 声学支对热容量的贡献}

\textit{模式密度}：声学支色散 $\omega=\omega_0'|\sin(qa)|$。一维波矢空间态密度为 $L/(2\pi)$，考虑 $\pm q$：
\begin{equation}
D(\omega)\,d\omega=2\cdot\frac{L}{2\pi}\,dq=\frac{Na}{\pi}\,dq.
\end{equation}

由 $d\omega/dq=\omega_0'a|\cos(qa)|=a\sqrt{\omega_0'^2-\omega^2}$，得
\begin{equation}
D(\omega)=\frac{2N}{\pi}\frac{1}{\sqrt{\omega_0'^2-\omega^2}},\qquad 0\le\omega\le\omega_0'.
\label{eq:1d-Domega}
\end{equation}
（验证：$\int_0^{\omega_0'}D(\omega)\,d\omega=N$，恰为 $N$ 个声学模。）

\textit{热容量}：声学支总振动能
\begin{equation}
E_-=\int_0^{\omega_0'}\frac{\hbar\omega}{e^{\hbar\omega/k_BT}-1}\,D(\omega)\,d\omega.
\end{equation}

故
\begin{equation}
C_V^-=k_B\int_0^{\omega_0'}\Bigl(\frac{\hbar\omega}{k_BT}\Bigr)^2\frac{e^{\hbar\omega/k_BT}}{\bigl(e^{\hbar\omega/k_BT}-1\bigr)^2}\,D(\omega)\,d\omega.
\end{equation}

代入 \eqref{eq:1d-Domega}：
\begin{equation}
C_V^-=\frac{2Nk_B}{\pi}\int_0^{\omega_0'}\frac{1}{\sqrt{\omega_0'^2-\omega^2}}\Bigl(\frac{\hbar\omega}{k_BT}\Bigr)^2\frac{e^{\hbar\omega/k_BT}}{\bigl(e^{\hbar\omega/k_BT}-1\bigr)^2}\,d\omega.
\end{equation}

\textbf{Step 3: 高温极限 ($T\gg\Theta_E,\Theta_D$)}

\textit{光学支}：$e^{\Theta_E/T}-1\approx\Theta_E/T$，
\begin{equation}
C_V^+\approx Nk_B\Bigl(\frac{\Theta_E}{T}\Bigr)^2\frac{1}{(\Theta_E/T)^2}=Nk_B.
\end{equation}

\textit{声学支}：$e^{\hbar\omega/k_BT}-1\approx\hbar\omega/k_BT$，被积函数化简为
\begin{equation}
\Bigl(\frac{\hbar\omega}{k_BT}\Bigr)^2\frac{1}{(\hbar\omega/k_BT)^2}\cdot\frac{1}{\sqrt{\omega_0'^2-\omega^2}}=\frac{1}{\sqrt{\omega_0'^2-\omega^2}}.
\end{equation}

故
\begin{equation}
C_V^-\approx\frac{2Nk_B}{\pi}\int_0^{\omega_0'}\frac{d\omega}{\sqrt{\omega_0'^2-\omega^2}}
=\frac{2Nk_B}{\pi}\cdot\frac{\pi}{2}=Nk_B.
\end{equation}

\textbf{高温总热容量}：
\begin{equation}
C_V=C_V^++C_V^-\approx Nk_B+Nk_B=2Nk_B.
\end{equation}
符合杜隆-珀蒂定律：共 $2N$ 个原子，每个原子贡献 $k_B$。

\textbf{Step 4: 低温极限 ($T\ll\Theta_E,\Theta_D$)}

\textit{光学支}：$\Theta_E/T\gg1$，$e^{\Theta_E/T}\gg1$，
\begin{equation}
C_V^+\approx Nk_B\Bigl(\frac{\Theta_E}{T}\Bigr)^2e^{-\Theta_E/T}\to0.
\end{equation}
光学支被指数冻结，贡献可忽略。

\textit{声学支}：低温下只有 $\omega\ll\omega_0'$ 的低频模式被激发，可作近似：
\begin{itemize}
    \item 分母 $\sqrt{\omega_0'^2-\omega^2}\approx\omega_0'$（常数）
    \item 积分上限延拓至 $\infty$（因高频端被玻色因子截断）
    \item $e^{\hbar\omega/k_BT}-1\approx e^{\hbar\omega/k_BT}$（$\hbar\omega\gg k_BT$ 不成立，但可用精确形式）
\end{itemize}

令 $x=\hbar\omega/k_BT$，$d\omega=k_BT\,dx/\hbar$：
\begin{align}
C_V^-&\approx\frac{2Nk_B}{\pi\omega_0'}\int_0^{\infty}\Bigl(\frac{k_BTx}{\hbar}\Bigr)^2\frac{e^x}{(e^x-1)^2}\,\frac{k_BT}{\hbar}\,dx\cdot\frac{\hbar^2}{k_B^2T^2}\\
&=\frac{2Nk_B}{\pi\omega_0'}\cdot\frac{k_BT}{\hbar}\int_0^{\infty}\frac{x^2e^x}{(e^x-1)^2}\,dx.
\end{align}

利用 $\displaystyle\int_0^{\infty}\frac{x^2e^x}{(e^x-1)^2}\,dx=2\zeta(2)=\frac{\pi^2}{3}$，得
\begin{equation}
C_V^-\approx\frac{2Nk_B}{\pi\omega_0'}\cdot\frac{k_BT}{\hbar}\cdot\frac{\pi^2}{3}
=\frac{2\pi Nk_B^2T}{3\hbar\omega_0'}\propto T.
\end{equation}

\textbf{低温总热容量}：
\begin{equation}
C_V=C_V^++C_V^-\approx0+\frac{2\pi Nk_B^2T}{3\hbar\omega_0'}\propto T.
\end{equation}
\end{derivationbox}

\begin{formulabox}[答案汇总]
\begin{center}
\begin{tabular}{cccc}
\toprule
温度范围 & 光学支 $C_V^+$ & 声学支 $C_V^-$ & 总热容量 $C_V$ \\
\midrule
一般式 & $Nk_B(\Theta_E/T)^2\dfrac{e^{\Theta_E/T}}{(e^{\Theta_E/T}-1)^2}$ & $\dfrac{2Nk_B}{\pi}\displaystyle\int_0^{\omega_0'}\dfrac{x^2e^x}{(e^x-1)^2}\dfrac{d\omega}{\sqrt{\omega_0'^2-\omega^2}}$ & $C_V^++C_V^-$ \\[12pt]
高温 & $\approx Nk_B$ & $\approx Nk_B$ & $\approx2Nk_B$ \\[8pt]
低温 & $\sim e^{-\Theta_E/T}\to0$ & $\propto T$ & $\propto T$ \\
\bottomrule
\end{tabular}
\end{center}
\end{formulabox}

\begin{insightbox}[物理图像与概念深化]
\textbf{1. 为什么光学支能用爱因斯坦模型？}

$M\gg m$ 意味着重原子几乎不动，轻原子在``固定''的势能阱中振动。光学支描述的是轻原子相对于重原子的反相振动，其频率主要由轻原子质量和恢复力常数决定，与相邻原胞的相位关系（$q$）很弱——就像一个被弹簧连在墙上的小球，弹簧另一端固定。这就是爱因斯坦模型的物理：$N$ 个独立的、频率相同的谐振子。

\textbf{2. 为什么声学支等效为单原子链？}

声学支描述的是同相振动（长波极限下两个子晶格一起动）。由于 $M\gg m$，轻原子 $m$ 像``寄生虫''一样跟着重原子 $M$ 一起运动，其惯性贡献可忽略。有效振动质量就是 $M$，晶格常数就是 $2a$（同种原子间距），于是退化为单原子链。

\textbf{3. 一维模式密度的范霍夫奇点}

\eqref{eq:1d-Domega} 中 $D(\omega)\propto1/\sqrt{\omega_0'^2-\omega^2}$ 在 $\omega=\omega_0'$ 处发散，这是\textbf{一维范霍夫奇点}（Van Hove singularity）的典型特征。来源于一维色散在布里渊区边界处群速度 $d\omega/dq\to0$（能带顶端平坦），态密度堆积导致发散。奇点本身可积，不影响热容量的有限性。

\textbf{4. 为什么低温下 $C_V\propto T$ 而非 $T^3$？}

\begin{center}
\begin{tabular}{ccc}
\toprule
维度 & 模式密度 & 低温热容 \\
\midrule
一维 & $D(\omega)\propto\omega^0$（常数） & $C_V\propto T$ \\
二维 & $D(\omega)\propto\omega$ & $C_V\propto T^2$ \\
三维 & $D(\omega)\propto\omega^2$ & $C_V\propto T^3$ \\
\bottomrule
\end{tabular}
\end{center}

低温下有效频率范围 $\omega_{\text{cutoff}}\sim k_BT/\hbar$。一维中 $D(\omega)$ 在此区间积分得 $N_{\text{modes}}\propto T^1$（而非三维的 $T^3$），每个模式能量 $\sim k_BT$，故 $U\propto T^2$，$C_V=dU/dT\propto T$。

\textbf{5. 杜隆-珀蒂定律的适用条件}

高温 $C_V\to2Nk_B$ 意味着每个振动自由度贡献 $k_B$。一维双原子链每个原胞有 2 个原子，每个原子 1 个自由度，共 $2N$ 个简正模，故总热容为 $2Nk_B$。这与经典能量均分定理一致：每个谐振子平均能量 $k_BT$（动能 $k_BT/2$ + 势能 $k_BT/2$）。
\end{insightbox}

\begin{thinkbox}[考场速记：一维双原子链热容问题的四步法]
遇到``一维双原子链求热容''类问题：
\begin{enumerate}
    \item \textbf{写色散并取近似}：$M\gg m$ 时光学支 $\omega_+\approx\text{const}$（爱因斯坦），声学支 $\omega_-\approx\omega_0'|\sin(qa)|$（德拜/单原子链）
    \item \textbf{分别求两支热容}：光学支直接爱因斯坦公式；声学支先求 $D(\omega)$ 再积分
    \item \textbf{高温展开}：$e^x-1\approx x$，验证杜隆-珀蒂 $C_V\to2Nk_B$
    \item \textbf{低温近似}：光学支指数冻结；声学支低频近似 $D(\omega)\approx\text{const}$，得 $C_V\propto T$
\end{enumerate}
\textit{核心记忆}：$M\gg m$ = 光学支（轻原子独立振动）+ 声学支（重原子单原子链）。一维低温热容永远是 $\propto T$，不是 $T^3$。
\end{thinkbox}
